\documentclass[aps,prl,amsmath,twocolumn,floatfix,superscriptaddress,longbibliography,footinbib]{revtex4-1}

\usepackage{amsmath, amsfonts, amssymb, bm, bbm}
\usepackage{mathtools}
\usepackage[inline]{enumitem}
\usepackage{graphicx}
\usepackage{xcolor}
\usepackage[breaklinks]{hyperref}
%\hypersetup{colorlinks,linkcolor={magenta},citecolor={blue},urlcolor={blue}} 
\hypersetup{
	pdfnewwindow=true,      % links in new window
	colorlinks=true,       % false: boxed links; true: colored links
	linkcolor=blue,          % color of internal links
	citecolor=magenta,        % color of links to 
	filecolor=blue,      % color of file links
	urlcolor=blue        % color of external links
}

% ORCID
\usepackage{sidecap,tikz}
\definecolor{lime}{HTML}{A6CE39}
\DeclareRobustCommand{\orcidicon}{\hspace{-1mm}
	\begin{tikzpicture}
		\draw[lime, fill=lime] (0,0) 
		circle [radius=0.16] 
		node[white] {{\fontfamily{qag}\selectfont \tiny \,ID}};
		\draw[white, fill=white] (-0.0525,0.095) 
		circle [radius=0.007];
	\end{tikzpicture}
	\hspace{-3mm}
}
\foreach \x in {A, ..., Z}{\expandafter\xdef\csname orcid\x\endcsname{\noexpand\href{https://orcid.org/\csname orcidauthor\x\endcsname}
		{\noexpand\orcidicon}}
}

\usepackage[capitalize]{cleveref} 

\usepackage{glossaries}
\glsdisablehyper 
\newacronym{wtd}{WTD}{waiting time distribution}
\newacronym{mbs}{MBS}{Majorana bound state}
\newacronym{qshi}{QSHI}{quantum spin Hall insulator}

% shortcuts
\newcommand{\up}{\uparrow}
\newcommand{\dw}{\downarrow}
\newcommand{\e}{\mathrm{e}}

\DeclarePairedDelimiter\mean{\langle}{\rangle}
\DeclarePairedDelimiter\abs{\lvert}{\rvert}
\DeclareMathOperator{\sgn}{sgn}
\DeclareMathOperator{\real}{Re}
\DeclareMathOperator{\imag}{Im}
\DeclareMathOperator{\diag}{diag}

\DeclareMathOperator{\sinc}{sinc}

\newcommand{\bea}{\begin{eqnarray}}
	\newcommand{\eea}{\end{eqnarray}}
\newcommand{\beq}{\begin{equation}}  
	\newcommand{\eeq}{\end{equation}}
\newcommand{\non}{\nonumber}

\newcommand\mb{\mathbf}
\newcommand\bs{\boldsymbol}
\newcommand\mT{\mathcal{T}}
\newcommand\dd{\text{d}}


\newcommand{\orcidauthorA}{0000-0003-2590-6231}
\newcommand{\orcidauthorB}{0000-0001-6037-6243}
\newcommand{\orcidauthorC}{0000-0002-4726-5247}
\newcommand{\orcidauthorD}{0000-0001-5726-0485}

\newcommand{\red}[1]{\textcolor{red}{#1}}
\newcommand{\lars}[1]{\textcolor{Mulberry}{{#1}}}
\newcommand{\green}[1]{\textcolor{green}{#1}}
\newcommand{\clr}[1]{\textcolor{blue}{#1}}

\begin{document} 

\title{Topological phase transition revealed by electron waiting times}

\author{Paramita Dutta\orcidA{}}
\affiliation{Theoretical Physics Division, Physical Research Laboratory, Ahmedabad-380009, India}

\author{Jorge Cayao\orcidB{}}
\affiliation{Department of Physics and Astronomy, Uppsala University, Box 516, S-751 20 Uppsala, Sweden}

\author{Annica M. Black-Schaffer\orcidC{}}
\affiliation{Department of Physics and Astronomy, Uppsala University, Box 516, S-751 20 Uppsala, Sweden}

\author{Pablo Burset\orcidD{}}
\affiliation{Department of Theoretical Condensed Matter Physics, Condensed Matter Physics Center (IFIMAC) and Instituto Nicol\'as Cabrera, Universidad Aut\'onoma de Madrid, 28049 Madrid, Spain}

\date \today

\begin{abstract}
The analysis of waiting times of electron transfers has recently become experimentally accessible owing to advances in noninvasive probes working in the short-time regime. We study  electron waiting times in a topological Andreev interferometer: a superconducting loop with controllable phase difference connected to a quantum spin Hall edge. The edge state helicity enables the transfer of electrons and holes into separate leads, with transmission controlled by the loop's phase difference $\phi$. In this setup, a topological phase transition with emerging Majorana bound states occurs at $\phi=\pi$. The waiting times for electron transfers across the junction are sensitive to the phase transition, but are uncorrelated for all $\phi$. By contrast, in the topological phase, the waiting times of Andreev-scattered holes show a strong correlation and the crossed (hole-electron) distributions feature a unique behavior. Both effects exclusively result from the presence of Majorana bound states. Consequently, electron waiting times could circumvent some of the challenges for detecting topological superconductivity and Majorana states beyond conductance signatures. 
\end{abstract}

\maketitle

Fluctuations in electron transport can greatly impact the performance of electronic circuits but, at the same time, provide us with invaluable information about the quantum-coherent behavior of conductors~\cite{Blanter2000}. 
Charge fluctuations are often analyzed by full counting statistics~\cite{Levitov1996,Bagrets2003,Bagrets2006}, which usually concerns the zero-frequency or long-time limit. However, experimental advances with noninvasive probes~\cite{Camus2017,Geller2019,Pekola2019} have now enabled access to the short-time regime with almost single-event resolution. 

A prominent, recently developed tool for the short-time regime is the electron waiting time distribution (WTD): the distribution of time intervals between consecutive charge transfers. Electron WTDs provide information about tunneling events in mesoscopic conductors, beyond average current and noise~\cite{Albert2011,Albert2012,Haack2014,Tang2014,Thomas2014,Sothman2014,Dasenbrook2015,Burset2019,Kosov2019,Kosov2021}, and have recently been measured in quantum dots connected to metallic electrodes~\cite{Gorman2017,Geller2019,Tarucha2020,Brange2021} and superconductors~\cite{Jenei2019,Ranni2021}.  
Moreover, correlations between waiting times indicate nonrenewal quantum transport~\cite{Kosov2019}, where electron transfers are not independent or identically distributed, making WTDs a source of information distinct from other statistical tools~\cite{Dasenbrook2015,Kosov2019,Kosov2021}. 


 %%%%%%%%%%%%%%%%%%%% FIGURE OF BAND STRUCTURE %%%%%%%%%%%
\begin{figure}[hb!]
	\centering
	\includegraphics[scale=0.65]{fig1a.pdf}
%	\hspace{-0.7cm}
	\includegraphics[scale=0.16]{fig1b.pdf}
	\caption{(a) Andreev interferometer at the QSHI edge with magnetic flux $\phi$ applied through a superconductor loop. (b) Incident electrons with energy within the gap $\Delta$ (blue balls) scatter off the NL-SL-SR-NR junction only as electron transmissions to NR or Andreev-reflected holes (red balls) into NL. Detectors D$_1$ and D$_2$, respectively placed at NL and NR, detect electrons or holes either individually or simultaneously. 
	NL is biased by a voltage $V$ and NR is at equilibrium. }
	\label{fig:model}
\end{figure}
%%%%%%%%%%%%%%%%%%%%%%%%%%%%%%%%%%%%%%%%%%%%%%%%%  

%In this Letter, we analyze the potential of electron WTDs and their correlations for revealing topological phase transitions, specifically, in topological superconductors hosting Majorana bound states (MBSs)~\cite{Kane2005,Fu2008,Fu2009,Qi2011}. There is currently an intense research activity focused on obtaining reliable signatures of MBSs, since the simplest one, a robust and quantized zero-bias conductance peak~\cite{Mourik2012,Higginbotham2015,Deng2016}, has proven insufficient~\cite{Frolov2019,Frolov2021}. In fact, electron waiting times in superconducting hybrid junctions have already been proposed to detect the presence of MBSs~\cite{Chevallier2016,Mi2018,Weymann2020,Fu2022,Schulz2022}, and also to characterize the entanglement between the electrons forming the Cooper pairs~\cite{Rajabi2013,Albert2016,Walldorf2018}. These theoretical proposals extended the concept of waiting times to deal with both the electron spin and the presence of holes, necessary for topological superconductors. However, to simplify the analysis, most approaches include ferromagnetic leads to reduce some degrees of freedom~\cite{Chevallier2016,Weymann2020,Schulz2022}. By contrast, our proposal does not involve any magnetic materials. 

In this Letter, we analyze the potential of electron WTDs and their correlations for identifying topological superconductors hosting Majorana bound states (MBSs)~\cite{Kane2005,Fu2008,Fu2009,Qi2011}. 
There is currently an intense research activity focused on obtaining reliable signatures of MBSs, since the simplest one, a robust and quantized zero-bias conductance peak~\cite{Mourik2012,Higginbotham2015,Deng2016}, has proven insufficient~\cite{Frolov2019,Frolov2021}. 
Electron waiting times in superconducting hybrid junctions have already been proposed to detect the presence of MBSs~\cite{Chevallier2016,Mi2018,Weymann2020,Fu2022,Schulz2022} and to characterize the entanglement between the electrons forming the Cooper pairs~\cite{Rajabi2013,Albert2016,Walldorf2018}. 
These theoretical proposals extended the concept of waiting times to deal with both the spin and electron-hole degrees of freedom. 
%However, to simplify the analysis, most approaches include ferromagnetic leads to reduce some degrees of freedom~\cite{Chevallier2016,Weymann2020,Schulz2022}. By contrast, our proposal does not involve any magnetic materials. 
In contrast to most approaches that include ferromagnetic leads to simplify the analysis~\cite{Chevallier2016,Weymann2020,Schulz2022}, our proposal does not involve any magnetic materials. 
%
%\clr{Due} to the complexities involved in measuring electron waiting times and engineering topological superconductors, we here suggest a Majorana platform that is both conceptually simple and presents important advantages for measuring waiting times of electrons and holes: 
%
Instead, we here suggest a Majorana platform that is both conceptually simple and presents important advantages for measuring waiting times of electrons and holes: 
an Andreev interferometer built on the edge of a quantum spin Hall insulator (QSHI)~\cite{Konig2007,Brune2010,Sullivan2011,Brune2012,Reis2017,Kammhuber2017,Wu2018} [\cref{fig:model}(a)]. 
%As we now detail, the main advantages of our proposal are the absence of magnetic materials and a unique way of separating electrons and holes into different leads that facilitates both the measurement of their waiting times and the study of the nonlocal nature of Majorana states.

The QSHI features helical edge states consisting of one-dimensional Dirac fermions characterized by spin-momentum locking~\cite{Wu2006,Xu2006}. 
When proximitized by a narrow superconducting lead~\cite{Sullivan2012,Hart2014,Bocquillon2016,Bocquillon2017,Deacon2017,Hart2017,Sajadi2018,Fatemi2018}, the helical edge states guarantee that only electrons tunnel through the lead and only Andreev-converted holes are reflected~\cite{Adroguer2010,Black-Schaffer2012,Black-Schaffer2013,Crepin2015,Tkachov2015,Cayao2017,Keidel2018,Keidel2020,Cayao2022,Lu2022}. 
Consequently, the superconductor acts as a beam splitter that separates electrons from holes into different leads and allows us to detect them independently, see \cref{fig:model}(b). 
%Moreover, a topological phase transition where MBSs are created can be controlled on-demand in an Andreev interferometer circuit~\cite{Chandrasekhar1998,Chandrasekhar2001}. 
In the interferometer setup~\cite{Chandrasekhar1998,Chandrasekhar2001}, a superconducting loop with controllable phase difference $\phi$ is connected to the QSHI edge [\cref{fig:model}(a)], so that
%Depending on the separation between the SL and SR loop arms, one or more Andreev bound states will form, facilitating transport across the whole SL-SR junction. As a result, 
the electric conductance at the NL or NR sides depends periodically on $\phi$. 
%In addition, when the phase difference is set to 
Interestingly, a topological phase transition occurs at $\phi=\pi$ where a pair of MBSs form at the SL-SR interface~\cite{Fu2008,Fu2009}. 
The MBS emergence can thus be controlled on-demand in the Andreev interferometer circuit. 

%We here analyze the electron waiting times and their correlations in the Andreev interferometer in \cref{fig:model} using a multi-terminal setup, which accounts for the independent and even simultaneous detection of electrons and holes, and especially compare the trivial and topological phases. 
%
We find that the waiting times for electron transfers across 
%the SL-SR 
our junction are sensitive to the formation of MBSs, but are still completely uncorrelated. 
%
By contrast, the waiting times of Andreev reflected holes are less sensitive to the MBSs, but instead present a strong correlation in the topological phase. 
%
Also importantly, the crossed (hole-electron) distributions and their correlations feature a unique behavior in the topological phase. 
%
% Measuring WTDs without QDs
Due to the sensitivity to the topological phase transition, electron waiting times and their correlations are an alternative signature of MBSs, circumventing the problems arising from resonant levels that naturally form in many Majorana platforms even in the trivial phase~\cite{Frolov2019}. 

%-----------------------------------------------
\textit{Topological Andreev interferometer.---}
%We consider an Andreev interferometer at the edge of a QSHI, which comprises of a superconducting loop with a short constriction between the SL and SR arms, in total forming the junction that is attached to the normal metal leads, NL and NR, see \cref{fig:model}. 
We consider an Andreev interferometer at the edge of a QSHI (\cref{fig:model}), which comprises of a superconducting loop with a short SL-SR junction that is attached to the normal metal leads NL and NR. 
For simplicity, we fix the length of each superconductor segment to be equal, $L_{\text S}$, and only analyze the situation where SL and SR share an interface~[\cref{fig:model}(b)]~\footnotemark[1]\footnotetext{A finite separation of the SL and SR arms, or asymmetry in the superconductor lengths, would not qualitatively affect our results.}. Low-energy excitations are described in the basis $\Psi(x)=(
\psi_{\uparrow},
\psi_{\downarrow},
\psi_{\downarrow}^{\dagger},
-\psi_{\uparrow}^{\dagger})^{T}$, with $\psi^{\dagger}_{\sigma}(x)$ the creation operator for electrons with spin $\sigma\in\{\up,\dw\}$ at position $x$, by the Bogoliubov-de Gennes Hamiltonian~\cite{Cayao2017} 
\begin{equation}\label{Eq:BdG_H}
H_{\rm BdG}=\hbar v_{F}k_{x} \hat{\eta}_{3}\hat{\sigma}_{3}-\mu \hat{\eta}_{3}\hat{\sigma}_{0} + \Delta (x) \hat{\eta}_{1} \hat{\sigma}_{0} .
\end{equation}
Here, $v_{F}$ is the Fermi velocity, $\mu$ the chemical potential, and the Pauli matrices $\hat{\eta}_{j}$ and $\hat{\sigma}_{j}$ act in Nambu and spin spaces, respectively. 
We set the pair potential $\Delta(x)\!=\!\Delta$ for SL, $\Delta(x)\!=\!\Delta e^{i \phi}$ for SR ($\phi$ is the superconducting phase difference), and zero otherwise. 
Henceforth, we set $v_F\!=\!1$, $\hbar\!=\!1$, and $\Delta\!=\!1$ so that the superconducting coherence length is $\xi\!=\!\hbar v_F/\Delta\!=\!1$~\footnotemark[2].
\footnotetext{Superconductor-normal metal-superconductor junctions with the normal intermediate region much smaller than the superconducting coherence length are usually referred to as \textit{short} junctions. We work in this regime, but use the terms \textit{short} and \textit{long} junctions to distinguish the cases where the total size of the SL-SR segment ($2L_\text{S}$) is smaller or larger, respectively, than the superconducting coherence length $\xi$.}
We show all results for $\mu\!=\!0$ and $eV=\Delta/2$~\footnotemark[3],
\footnotetext{Our results are mostly insensitive to the chemical potential and do not qualitatively change for voltages within the gap.} 
see Supplementary Material (SM)~\cite{SM}. 
The QSHI Andreev interferometer forms a topological Josephson-like junction that hosts MBSs at the SL-SR interface when $\phi=\pi$~\cite{Fu2009,Cayao2022}. This \emph{topological phase transition} separates two distinct phases: with and without MBSs in the middle of the junction. Subsequently, we compare the cases with $\phi=0$ and $\phi=\pi$, which we term, respectively, as the \textit{trivial} and \textit{nontrivial} phases. 

\textit{Electron waiting times.---}
WTDs for phase-coherent transport of noninteracting electrons are evaluated from the scattering matrix~\cite{Albert2012,Haack2014,Mi2018}. 
Generally an Andreev interferometer has four effective transport channels [electrons or holes (e, h), incoming or outgoing (i, o), from the left or right leads (L, R)], represented by the spinor $\Psi^{i(o)}=(\psi_{eL}^{i,(o)}, \psi_{hL}^{i(o)}, \psi_{eR}^{i(o)}, \psi_{hR}^{i(o)})^T$. For a given energy $E$, the scattering matrix connects outgoing and incoming solutions of \cref{Eq:BdG_H} as $\Psi^{(o)}= \mathcal{S}\Psi^{(i)}$~\cite{SM}. 
Owing to the spin-momentum locking at the QSHI edge, here only the normal transmissions $\mathcal{S}_{\alpha L,\alpha R}$ and $ \mathcal{S}_{\alpha R,\alpha L}$, with $\alpha=e,h$, and the Andreev reflections $ \mathcal{S}_{e X, hX}$ and $ \mathcal{S}_{h X, e X}$, with $X=L,R$, are nonzero. 
From the scattering matrix we define the idle-time probability $\Pi(\{\tau_\gamma\})$ that no particles of type $\gamma=\alpha X$ are detected during the time interval $\tau_\gamma$ (for the stationary processes considered here only time intervals are relevant). Following Refs.~\cite{Haack2014, Mi2018}, we have
\begin{equation}\label{eq:itp}
	\Pi(\{\tau_\gamma\})= \mathrm{det}[\mathcal{I}-\mathcal{S}^{\dagger}(E) \mathcal{K}(\{\tau_\gamma\},E-E') \mathcal{S}(E') ] , 
\end{equation}
where $\mathcal{I}$ is the identity matrix and $\mathcal{K}$ is a diagonal matrix, 
$\mathcal{K}(\{\tau_\gamma\},E) =\bigoplus_{\gamma} K( \tau_\gamma, E)$,
with the kernels~\cite{SM}
\begin{equation}\label{eq:kernel}
	K(\tau_\gamma, E)= \kappa \e^{-i E \tau_\gamma/2 } \sin (E \tau_\gamma/2) / (\pi E ) . 
\end{equation}
The linear dispersion relation of the QSHI helical edge states allows us to naturally divide the transport window $[\mu , \mu + eV ]$ in intervals of width $\kappa=eV/N$, where $N$ is the total number of intervals and $eV$ the applied bias. Due to the inversion symmetry of our setup, we only consider voltages applied to the left lead, $V_L\equiv V$, $V_R=0$. We work in the limit $N\rightarrow\infty$ where \cref{eq:kernel} correctly applies to stationary transport~\cite{Haack2014}. 

We can now define $\mathcal{W}_{\alpha\beta}(\tau)= -\mean{\tau_\alpha}\partial_\tau^2 \Pi(\tau)$ as the probability density of detecting a particle of type $\beta$ at a time $\tau$ after having measured a particle of type $\alpha$~\cite{SM}. Here, the mean waiting time $\mean{\tau_\alpha}$ is related to the average current for $\alpha$ particles, $I_\alpha=1/\mean{\tau_\alpha}$. 
Analogously, we define the joint waiting time $\mathcal{W}_{\alpha\gamma\beta}(\tau_1,\tau_2)= \mean{\tau_\alpha} \partial_{\tau_1} \partial_{\tau_2}^2\Pi$, which generalizes the waiting time distribution between particles of type $\alpha$ and $\beta$ to include the extra detection of a particle of type $\gamma$ at an intermediate time $\tau_1$, such that $0 \leq \tau_1 \leq \tau_2$~\cite{SM}. 
% Correlations
The joint WTD describes correlations between consecutive waiting times. When the waiting times are uncorrelated, the joint distribution factorizes as the product of two waiting time distributions~\cite{Dasenbrook2015}, $\mathcal{W}^\text{unc}_{\alpha\gamma\beta}(\tau_1,\tau_2)= \mathcal{W}_{\alpha\gamma}(\tau_1)\mathcal{W}_{\gamma\beta}(\tau_2)$. We can further quantify the correlations between consecutive waiting times using the correlation function
\begin{equation}\label{eq:correlations}
	\delta\mathcal{W}_{\alpha\gamma\beta}(\tau_1,\tau_2) = \frac{\mathcal{W}_{\alpha\gamma\beta}(\tau_1,\tau_2) - \mathcal{W}_{\alpha\gamma}(\tau_1)\mathcal{W}_{\gamma\beta}(\tau_2)}{\mathcal{W}_{\alpha\gamma}(\tau_1)\mathcal{W}_{\gamma\beta}(\tau_2)} .
\end{equation}

A main feature of the QSHI topological Andreev interferometer is that for an electron (say, spin-up) injected in NL, only (spin-down) holes and (spin-up) electrons can scatter into NL and NR, respectively. Thus, with electrons and holes always scattering into different leads, all \textit{local} or same detector WTDs are necessarily given by $\mathcal{W}_{ee}$ and $\mathcal{W}_{hh}$, while all \textit{nonlocal} WTDs are given by $\mathcal{W}_{eh}$ and $\mathcal{W}_{he}$, where measurements take place at different detectors. Similarly, $\mathcal{W}_{eee}$ and $\mathcal{W}_{hhh}$ are local joint WTDs, while joint distributions combining electron (NR) and hole (NL) measurements, like $\mathcal{W}_{ehe}$, are nonlocal.
%%%%%%%%%%%%%%%%%%%%%%%%%%%%%%%
\begin{figure}[t!]
	\centering
	\includegraphics[scale=0.5]{fig2a.pdf}
	\includegraphics[scale=0.49]{fig2b.pdf}
	\caption{Local WTDs, $\mathcal{W}_{ee}$ (a) and $\mathcal{W}_{hh}$ (b), for $L_{\rm S}=0.2$ (short junction, green) and $L_{\rm S}=1.5$ (long junction, magenta), with $\phi\! =\!0$ (trivial phase) for dashed lines and $\phi\! =\!\pi$ (topological phase) for solid ones. 
	The gray (orange) shaded region indicate the Poisson (Wigner-Dyson) distribution. }
	\label{fig:localWTD}
\end{figure}
%%%%%%%%%%%%%%%%%%%%%%%%%%%%%%%%


\textit{Local waiting times.---} 
We find that the local WTDs $\mathcal{W}_{\alpha\alpha}$, representing two consecutive detections at either NL ($\alpha=h$) or NR ($\alpha=e$), follow a behavior similar to that of normal-superconductor hybrid junctions in the trivial phase~\cite{Chevallier2016,Mi2018,Schulz2022}. It was established in Ref.\,[\onlinecite{Haack2014}] that the WTD of a quantum-coherent channel with energy-independent transmission is determined by its scattering probability: highly-transmitting channels result in a Wigner-Dyson distribution [orange area in \cref{fig:localWTD}(a)], describing a coherent particle flow, while low transparent channels result in a Poisson distribution [gray area in \cref{fig:localWTD}(a)], characteristic of tunnel transport. In both cases, the impossibility of a simultaneous measurement of two particles at the same detector due to the Pauli exclusion principle forces the WTDs to be zero at $\tau=0$. 
Owing to the constrained transport at the QSHI edge, with, e.g., $|\mathcal{S}_{hL,eL}|^2+|\mathcal{S}_{eR,eL}|^2=1$, the transition between the Wigner-Dyson and Poisson distributions is controlled by the length of the Andreev interferometer, $2L_\text{S}$. Long (short) junctions with $L_\text{S}>\xi$ ($L_\text{S}<\xi$) give a high probability of Andreev reflection, $|\mathcal{S}_{hL,eL}|^2$ (electron transmission, $|\mathcal{S}_{eR,eL}|^2$), and result in a Poisson (Wigner-Dyson) distribution~\cite{SM}. For the trivial phase $\phi=0$, represented in \cref{fig:localWTD} by dashed lines, we find these features for long and short junctions~\footnotemark[2], respectively, which is consistent with earlier results~\cite{Haack2014,Mi2018,Schulz2022} since the scattering probabilities for the topological Andreev interferometer are almost constant at subgap energies~\cite{SM}. 

By contrast, in the nontrivial phase found at $\phi=\pi$, the low-energy electron transmission probability becomes strongly energy-dependent for long junctions due to the resonant-tunneling through the MBS~\cite{SM,Cayao2022}. Consequently, we find that $\mathcal{W}_{ee}$ converges to the WTD of a resonant level in the tunnel limit~\cite{Haack2014}, see solid magenta line in \cref{fig:localWTD}(a), instead of evolving into a Poisson distribution like for $\phi=0$. The high-transmission (short junction) limit is less affected by the presence of the MBS, so we see little change with $\phi$. The distribution for reflected holes, $\mathcal{W}_{hh}$, which followed the opposite behavior to $\mathcal{W}_{ee}$ in the trivial phase with respect to junction length, is overall less affected by the presence of the MBS, see \cref{fig:localWTD}(b). This behavior we attribute to the continually high probability of Andreev reflection ($|\mathcal{S}_{hL,eL}|^2\sim1$). 

%%%%%%%%%%%%%%%%%%%%%%%%%%%%%%%%%%%
\begin{figure}[t!]
	\centering
	\includegraphics[scale=0.32]{fig3a.png}
	\includegraphics[scale=0.32]{fig3b.png}
	\caption{Correlations between waiting times for three consecutive hole detections at NL, $\delta\mathcal{W}_{hhh}(\tau_1,\tau_2)$, for $L_{\rm S}\!=\!1.5$ as a function of $\tau_1$ and $\tau_2$ at (a) $\phi\!=\!0$ and (b) $\phi\!=\!\pi$. }
	\label{fig:jointWTD}
\end{figure}
%%%%%%%%%%%%%%%%%%%%%%%%%%%%%%%%%%%
 
Even though $\mathcal{W}_{hh}$ is not very sensitive to the topological phase transition, the correlations between consecutive waiting times for hole transfers, $\delta\mathcal{W}_{hhh}$, contain very relevant information. We focus on long junctions, $L_\text{S}>\xi$, which are dominated by Andreev reflection processes and feature a more distinct behavior between the topological and trivial phases. In the trivial phase, the Andreev interferometer behaves like an electron-hole beam splitter, featuring the same correlations as a standard quantum point contact for electrons~\cite{Dasenbrook2015}, see \cref{fig:jointWTD}(a).
When the time between two hole transfers is small, $\tau_1<\mean{\tau_h}$ (or long, $\tau_1>\mean{\tau_h}$), the next hole detection at $\tau_2$ will require a long (short) waiting time (red color signals positive correlations). 
By contrast, in the topological phase we find correlations that seemingly explode at short waiting times, with values increasing an order of magnitude compared to the trivial case, see \cref{fig:jointWTD}(b). This means that short time intervals between detections are the most likely. 
The waiting times between electron transfers, on the other hand, seem to be completely uncorrelated, i.e., $\delta\mathcal{W}_{eee}(\tau_1,\tau_2)\simeq \mathcal{W}_{ee}(\tau_1)\mathcal{W}_{ee}(\tau_2)$~\cite{SM}. 

 %%%%%%%%%%%%%%%%%%%%%%%%%%%%%%%
\begin{figure}[t!]
	\centering
	\includegraphics[scale=0.5]{fig4a.pdf}
	\includegraphics[scale=0.5]{fig4b.pdf}
	\caption{Nonlocal WTDs, $\mathcal{W}_{he}$ (a) and $\mathcal{W}_{eh}$ (b), for $L_{\rm S}\!=\!0.2$ (short junction, green) and $L_{\rm S}=1.5$ (long junction, magenta), with $\phi\! =\!0$ (trivial phase) for dashed lines and $\phi\! =\!\pi$ (topological phase) for solid ones. }
	\label{fig:nonlocalWTD}
\end{figure}
%%%%%%%%%%%%%%%%%%%%%%%%%%%%%%%

\textit{Nonlocal waiting times.---} 
We next fully exploit the multi-terminal advantage of the topological Andreev interferometer by exploring the nonlocal WTDs. By definition, the distribution $\mathcal{W}_{\alpha\beta}$, with $\beta\neq\alpha$, assumes that the first particle $\alpha$ has been detected; no matter how unlikely that event is. Therefore, the nonlocal waiting times are determined by the probability of the second detection. Consequently, $\mathcal{W}_{eh}$ ($\mathcal{W}_{he}$) is determined by the Andreev reflection (electron transmission) probability, following a behavior similar to $\mathcal{W}_{hh}$ ($\mathcal{W}_{ee}$), which we verify in \cref{fig:nonlocalWTD}. The one marked difference between local and nonlocal distributions is that, as particles transfer into different detectors, nonlocal WTDs can be finite at zero waiting time and also always fulfill $\mathcal{W}_{eh}(0)=\mathcal{W}_{he}(0)$~\cite{Mi2018}. 

The nonlocal WTDs at zero waiting time have already been established to increase in the topological phase, independently of the scattering probabilities and due to the emergence of MBSs~\cite{Mi2018}. Here, we interestingly also find that $\mathcal{W}_{he}(\tau)$, which for $\phi=\pi$ is determined by the Majorana-assisted electron tunneling, is further strongly altered. Specifically, $\mathcal{W}_{he}(\tau)$ presents a dip at short but finite waiting times. We explain this behavior as being due to the transition between a regime dominated by the Andreev reflection probability at $\tau\rightarrow0$ into a regime where electron transmissions dominate at long waiting times. The former initially reduces the probability, while the latter imposes a behavior similar to the local distribution, $\mathcal{W}_{ee}$. The dip, or local minimum at short waiting times, reflects this transition and is particularly visible in presence of Majorana-induced resonant tunneling when the corresponding local WTD becomes anomalous, see \cref{fig:localWTD}(a). With $\mathcal{W}_{eh}$ being primarily determined by $\mathcal{W}_{hh}$, we find no such dip in $\mathcal{W}_{eh}$. This behavior of $\mathcal{W}_{he}$ is unique to the topological Andreev interferometer, which we confirmed by checking both local and nonlocal WTDs for an ordinary interferometer, in the absence of any topology. 

We further find that the correlations between nonlocal waiting times also show a unique behavior at the topological phase. We focus on alternate electron-hole-electron transfers in the long junction regime, where the contrast between the topological and trivial phases is more pronounced, and study the behavior of both $\mathcal{W}_{he}$ and $\mathcal{W}_{ee}$. Very short waiting times, $\tau_1<\mean{\tau_{h}}$, show only a weak correlation with long waiting times ($\tau_2>\mean{\tau_{e}}$) in the trivial phase, but this behavior is enhanced two orders of magnitude in the topological phase, see \cref{fig:jointWTD2}. As mentioned above, the sequential tunneling of three electrons, $\mathcal{W}_{eee}$, is uncorrelated. However, including one hole detection between the electron transfers drastically changes the statistics in the topological phase. 
These results indicate that to completely characterize the topological phase in an Andreev interferometer, we need to compare transport processes on both arms of the circuit. 
We note that $\mathcal{W}_{heh}$, which is dominated by Andreev reflections and thus less sensitive to the presence of the SL-SR junction, shows negative correlations and a weak phase dependence~\cite{SM}. The behavior of all studied WTDs and correlations is summarized in \cref{WTDs_All}. 

%%%%%%%%%%%%%%%%%%%%%%%%%%%%%%%%%%%
\begin{figure}[t!]
	\centering
	\includegraphics[scale=0.331]{fig5a.png}
	\includegraphics[scale=0.327]{fig5b.png}
	\caption{Correlations between waiting times for the detection sequence electron-hole-electron between NL and NR: $\delta\mathcal{W}_{ehe}(\tau_1,\tau_2)$ for $L_{\rm S}\!=\!1.5$ as a function of $\tau_1$ and $\tau_2$ at (a) $\phi\!=\!0$ and (b) $\phi\!=\!\pi$. }
	\label{fig:jointWTD2}
\end{figure}
%%%%%%%%%%%%%%%%%%%%%%%%%%%%%%%%%%%

\textit{Concluding remarks.---} 
We have analyzed the distribution of waiting times for emitted electrons and holes from a topological Andreev interferometer: a NL-SL-SR-NR junction on the quantum spin Hall edge. Two special features of the topological Andreev interferometer are 
(i) the emergence of MBSs when the superconducting phase difference between SL and SR is $\phi=\pi$ and 
(ii) that it acts as an electron-hole beam splitter, sending holes and electrons to different leads. The topological Andreev interferometer is thus one of the simplest multi-terminal setups without magnetic elements that features MBSs and allows us to test their nonlocal behavior. We find that, while the waiting times between holes in the same lead are not sensitive to the phase transition, the transmitted electron and, most importantly, the \textit{crossed} (hole-electron) distributions, feature distinct behavior in the topological state, due to the MBSs.
Moreover, the correlations between waiting times are also very sensitive to the MBS in the topological phase.
%For example, the waiting times between hole detections are strongly correlated for $\phi=\pi$, even though their distribution is not very sensitive to the superconducting phase difference. 
We collect our main results in \cref{WTDs_All}. 

A few remarks are in order. Previous works have analyzed the WTD of electrons tunneling into a Majorana state~\cite{Chevallier2016,Mi2018,Fu2022}, but for notably different setups and with less distinctive results. 
Moreover, Ref.~[\onlinecite{Schulz2022}] appeared during the final stages of our manuscript preparation, also studying a setup based on the QSHI but including ferromagnetic regions, showing that the resonant transport through a MBS yields WTDs very similar to that of a single level resonance~\cite{Haack2014,Mi2018}. 
This is a common problem encountered in the tunnel spectroscopy of MBSs~\cite{Frolov2019,Frolov2021}. 
%, and showing that the resonant transport through a MBS yields a very specific type of WTD. 
%Such behavior is, however, very similar to that of a single level resonance~\cite{Haack2014,Mi2018}, like that of a quantum dot level. 
%%%%%%%%%%%%%%%%%%%%%%%%%%%%%%%%%%%%%%%%%%%%%%
\begingroup
\squeezetable
\begin{table}[t!]
	\begin{tabular}{|c|c|c|c|}
		\hline %\hline
		WTD &  $\phi\!=\!0$ &  $\phi\!=\!\pi$ \\ \hline 
		$\mathcal{W}_{ee}$ & P  & Resonant level \\  \hline 
		$\mathcal{W}_{hh}$ & WD  & WD \\ \hline 
		$\mathcal{W}_{eh}$ & $\mathcal{W}_{eh}(0)\approx 0$, WD & $\mathcal{W}_{eh}(0)> 0$, WD  \\ \hline 
		$\mathcal{W}_{he}$ & $\mathcal{W}_{he}(0)\approx 0$, P & $\mathcal{W}_{he}(0)> 0$, P with dip \\ \hline  \hline
		$\delta\mathcal{W}_{eee}$ & Uncorrelated  & Uncorrelated \\  \hline 
		$\delta\mathcal{W}_{hhh}$ & Bream-splitter & $\delta\mathcal{W}_{hhh}\sim 1$ at $\tau_{1,2}<\mean{\tau_h}$ \\  \hline 
		$\delta\mathcal{W}_{heh}$ & Low correlation  & Low correlation \\  \hline 
		$\delta\mathcal{W}_{ehe}$ & Low correlation & $\delta\mathcal{W}_{ehe}\sim0.1$ at $\tau_{1}<\mean{\tau_h}$ \\  \hline 
%		\hline
	\end{tabular}
	\caption{Summary of the behaviors of WTDs and joint WTDs for $L_{\rm S}\!>\!\xi$. Here, WD and P indicate, respectively, Wigner-Dyson and Poisson distributions. }
	\label{WTDs_All}
\end{table}
\endgroup
%%%%%%%%%%%%%%%%%%%%%%%%%%%%%%%%%%%%%%%%%%%%%%%
Our results complement these works by showing that taking into account independent and simultaneous transfers of electrons and holes into separate detectors (i.e., both local and nonlocal WTDs), and the correlations between them, yields a very distinctive signature of Majorana modes. 
Despite the challenges involved in measuring waiting times, there are promising new advances like experimental measurements of time-of-flight of electron excitations~\cite{Kataoka2016,Roussely2018}, in addition to recent theoretical proposals for WTD clocks~\cite{Haack2015,Dasenbrook2016}. 
%Although measuring electron waiting times in mesoscopic conductors remains an experimental challenge, there are promising new advances like experimental measurements of time-of-flight of electron excitations~\cite{Kataoka2016,Roussely2018}, in addition to recent theoretical proposals for WTD clocks~\cite{Haack2015,Dasenbrook2016}. 


 %-----------------------------------------------
\acknowledgments
P.D.~thanks D.~Chakraborty, P.~Holmvall, and I.~Mahyaeh for technical help, U.~Basu for discussions, and acknowledges the computational facilities by Uppsala University and the Department of Science and Technology (DST), India, for the financial support through SERB Start-up Research Grant (File No.~SRG/2022/001121). 
J.C. acknowledges financial support from the Swedish Research Council (Vetenskapsr{\aa}det Grant No. 2021-04121), the Scandinavia-Japan Sasakawa Foundation (Grant No. GA22-SWE-0028), the Royal Swedish Academy of Sciences (Grant No. PH2022-0003), and the Carl Trygger’s Foundation (Grant No. 22: 2093). 
A.B.S. acknowledges financial support from the Swedish Research Council (Vetenskapsr{\aa}det Grant No. 2018-03488) and the Knut and Alice Wallenberg Foundation through the Wallenberg Academy Fellows program and the EU-COST Action CA-21144 Superqumap. 
P.B.~acknowledges support from the Spanish CM ``Talento Program'' project No.~2019-T1/IND-14088 and the Agencia Estatal de Investigaci\'on projects No.~PID2020-117992GA-I00 and No.~CNS2022-135950. 

\bibliography{bibfile}

%\include{SMv3.tex}
%\end{document}

%\newpage

% Supplemental material
\clearpage

%
%
\onecolumngrid
%
%
\setcounter{equation}{0}
\renewcommand{\theequation}{S\,\arabic{equation}}
\setcounter{figure}{0}
\renewcommand{\thefigure}{S\,\arabic{figure}}
%
%
%-----------------------------------------------
%\begin{appendix}
%\appendix
%-----------------------------------------------
%\onecolumngrid
\section{Supplementary material for ``Topological phase transition revealed by electron waiting times"}

In this supplemental material, we provide details of the theoretical formalism and some additional results to support the discussions of the main text about the scattering probabilities, the bias dependence of the waiting time distributions, and the nonlocal joint waiting time correlations. 

\section{Theoretical methods} \label{sec:App-method}

In this section, we summarize the most important steps in order to arrive at the scattering matrices used in the main text and describe the theory for WTDs and joint WTDs.

\subsection{Scattering matrix formalism\label{apnd1}}

We consider a NL-SL-SR-NR junction, with NL and NR normal leads and SL and SR superconducting ones, along the x direction on one edge of a QSHI, where the NL-SL interface is placed at $x=0$, the SL-SR interface at $x=L_\text{S}$, and the SR-NR one at $x=2L_\text{S}$. The scattering states at the different regions of the NL-SL-SR-NR junction can be written as~\cite{Cayao2022}
\bea
\psi_{\rm 1}(x)&=&
 \phi_{1}^{N}\,{\rm e}^{ik_{e}x}+r_{\rm eh}^{\rm L}\phi_{3}^{N}\,{\rm e}^{ik_{h}x} +r_{\rm ee}^{\rm L}\phi_{2}^{N}\,{\rm e}^{-ik_{e}x},\,\quad x<0  \\
  &=&a_{1}\phi_{1}^{\rm S}\,{\rm e}^{ik^{\rm S}_{e}x}+b_{1}\phi_{2}^{\rm S}\,{\rm e}^{-ik^{\rm S}_{e}x}+c_{1}\phi_{3}^{\rm S}\,{\rm e}^{ik^{\rm S}_{h}x}+d_{1}\phi_{4}^{\rm S}\,{\rm e}^{-ik^{\rm S}_{h}x}\,,\quad 0<x<L_{\rm S} \non \\
&=&   p_{1}\phi_{1}^{\rm S}\,{\rm e}^{ik^{\rm S}_{e}x}+q_{1}\phi_{2}^{\rm S}\,{\rm e}^{-ik^{\rm S}_{e}x}+r_{1}\phi_{3}^{\rm S}\,{\rm e}^{ik^{\rm S}_{h}x}+s_{1}\phi_{4}^{\rm S}\,{\rm e}^{-ik^{\rm S}_{h}x}\,,\quad L_{\rm S}<x<2L_{\rm S} \non \\   
 &=& t_{\rm ee}^R  \phi_{1}^{\rm N}\,{\rm e}^{ik_{e}x} + t_{\rm eh}^{\rm R}  \phi_{4}^{\rm N}\,{\rm e}^{-ik_{h}x},\,x>2L_{\rm S} \non \\ 
\psi_{2}(x)&=&
 \phi_{4}^{N}\,{\rm e}^{-ik_{h}x}++r_{\rm he}^{\rm L}\phi_{2}^{N}\,{\rm e}^{-ik_{e}x} +r_{\rm hh}^{\rm L} \phi_{3}^{N}\,{\rm e}^{ik_{h}x},\,x<0 \\
 &=& a_{2}\phi_{1}^{\rm S}\,{\rm e}^{ik^{\rm S}_{e}x}+ b_{2}\phi_{2}^{\rm S}\,{\rm e}^{-ik^{\rm S}_{e}x}+ c_{2}\phi_{3}^{\rm S}\,{\rm e}^{ik^{\rm S}_{h}x}+ d_{2}\phi_{4}^{\rm S}\,{\rm e}^{-ik^{\rm S}_{h}x}\,,\quad 0<x<L_{\rm S} \non \\
  &=&p_{2}\phi_{1}^{\rm S}\,{\rm e}^{ik^{\rm S}_{e}x}+q_{2}\phi_{2}^{\rm S}\,{\rm e}^{-ik^{\rm S}_{e}x}+r_{2}\phi_{3}^{\rm S}\,{\rm e}^{ik^{\rm S}_{h}x}+s_{2}\phi_{4}^{\rm S}\,{\rm e}^{-ik^{\rm S}_{h}x}\,,\quad L_{\rm S}<x<2L_{\rm S} \non \\ 
   &=& t_{\rm hh}^{\rm R}  \phi_{4}^{\rm N}\,{\rm e}^{-ik_{h}x} + t_{\rm he}^{\rm R}  \phi_{1}^{\rm N}\,{\rm e}^{ik_{e}x}, \,x>2L_{\rm S} \non \\
\psi_{3}(x)&=& c_{3}  \phi_{2}^{N}\,{\rm e}^{-ik_{e}x} +d_{3}  \phi_{3}^{N}\,{\rm e}^{ik_{h}x},\, x<0  \\
    &=&a_{3}\phi_{1}^{\rm S}\,{\rm e}^{ik^{\rm S}_{e}x}+b_{3}\phi_{2}^{\rm S}\,{\rm e}^{-ik^{\rm S}_{e}x}+c_{3}\phi_{3}^{\rm S}\,{\rm e}^{ik^{\rm S}_{h}x}+d_{3}\phi_{4}^{\rm S}\,{\rm e}^{-ik^{\rm S}_{h}x}\,,\quad 0<x<L_{\rm S} \non \\
 &=& p_{3}\phi_{1}^{\rm S}\,{\rm e}^{ik^{\rm S}_{e}x}+q_{3}\phi_{2}^{\rm S}\,{\rm e}^{-ik^{\rm S}_{e}x}+r_{3}\phi_{3}^{\rm S}\,{\rm e}^{ik^{\rm S}_{h}x}+s_{3}\phi_{4}^{\rm S}\,{\rm e}^{-ik^{\rm S}_{h}x}\,,\quad L_{\rm S}<x<2L_{\rm S} \non \\ 
   &=& \phi_{2}^{\rm N}\,{\rm e}^{-ik_{e}x}+a_{3}\phi_{4}^{\rm N}\,{\rm e}^{-ik_{h}x}+b_{3}\phi_{1}^{\rm N}\,{\rm e}^{ik_{e}x},\,x>2L_{\rm S}\non \\
   \psi_{4}(x)&=& c_{4}  \phi_{3}^{N}\,{\rm e}^{ik_{h}x} +    d_{4}  \phi_{2}^{N}\,{\rm e}^{-ik_{e}x},\,x<0   \\
    &=&a_{4}\phi_{1}^{\rm S}\,{\rm e}^{ik^{\rm S}_{e}x}+ b_{4}\phi_{2}^{\rm S}\,{\rm e}^{-ik^{\rm S}_{e}x}+ c_{4}\phi_{3}^{\rm S}\,{\rm e}^{ik^{\rm S}_{h}x}+ d_{4}\phi_{4}^{\rm S}\,{\rm e}^{-ik^{\rm S}_{h}x}\,,\quad 0<x<L_{\rm S} \non \\
   &=&p_{4}\phi_{1}^{\rm S}\,{\rm e}^{ik^{\rm S}_{e}x}+ q_{4}\phi_{2}^{\rm S}\,{\rm e}^{-ik^{\rm S}_{e}x}+r_{4}\phi_{3}^{\rm S}\,{\rm e}^{ik^{\rm S}_{h}x}+s_{4}\phi_{4}^{\rm S}\,{\rm e}^{-ik^{\rm S}_{h}x}\,,\quad L_{\rm S}<x<2L_{\rm S} \non \\ 
&=&\phi_{3}^{\rm N}\,{\rm e}^{ik_{h}x}+a_{4}\phi_{1}^{\rm N}\,{\rm e}^{ik_{e}x}+b_{4}\phi_{4}^{\rm N}\,{\rm e}^{-ik_{h}x},\,x>2L_{\rm S} \non
\label{phiNS}
\eea
where
\beq
\begin{split}
\phi_{1}^{\rm N}\!&=\!\begin{pmatrix}
1,0,0,0
 \end{pmatrix}^{T},
\phi_{2}^{\rm N}\!=\!\begin{pmatrix}
0,1,0,0
 \end{pmatrix}^{T},
\phi_{3}^{\rm N}\!=\!\begin{pmatrix}
0,0,1,0
 \end{pmatrix}^{T},
\phi_{4}^{N}\!=\!\begin{pmatrix}
0,0,0,1
 \end{pmatrix}^{T},\\
 \phi_{1}^{\rm S}\!&=\!\begin{pmatrix}
u, 0,v,0
 \end{pmatrix}^{T},
\phi_{2}^{\rm S}\!=\!\begin{pmatrix}
0,u,0,v
 \end{pmatrix}^{T},
\phi_{3}^{\rm S}\!=\!\begin{pmatrix}
v,0,u,0
 \end{pmatrix}^{T},
\phi_{4}^{\rm S}\!=\!\begin{pmatrix}
0,v,0,u
 \end{pmatrix}^{T}.
\end{split}
\eeq
The $a_i$, $b_i$, $c_i$, and $d_i$ are the scattering parameters to be determined. The wave vectors (i) in normal regions NL and NR read as 
\bea
k_{e(h)}(E,\Delta)=\frac{\mu \pm E}{v_F},
\eea
and (ii) in the superconducting regions SL and SR take the form,
\bea
k_{e(h)}^S(E,\Delta)=\frac{\mu}{v_F}\pm\frac{\sqrt{E^2-\Delta^2}}{v_F} .
\nonumber
\eea
We solve the above equations by matching the wave functions at the interfaces $x=0$, $x=L_{\rm S}$, and $x=2L_{\rm S}$, to obtain 
\bea
b_{i}&=0\,,\quad
d_{i}=0\,,\quad
p_{2,3}=0\,,\quad
q_{1,4}=0\,,\quad
r_{2,3}=0\,, \non \\
s_{1,4}&=0\,,\quad
e_{1,4}=0\,,\quad
f_{2,3}=0\,,\quad
g_{1,4}=0\,,\quad
h_{2,3}=0. \non
\eea
The local normal reflection and nonlocal crossed Andreev transmission are forbidden by the conservation of helicity of the edge states of the QSHI~\cite{Adroguer2010,Tkachov2013} allowing only two processes: (i) local Andreev reflections, where an incident electron is reflected as a hole at NL, and (ii) nonlocal electron transmission at NR with the amplitudes as follows~\cite{Cayao2022}. 
\noindent
\begin{itemize}
\item Andreev reflection in NL for an incident electron from NL:
\beq\label{eq:ref-amp}
r^{\rm L}_{\rm eh}=\frac{(e^{i k_{\rm e}^{\rm S} L_{\rm S}} - e^{i k_{\rm h}^{\rm S} L_{\rm S}}) uv [(e^{i k_{\rm e}^{\rm S} L_{\rm S}} + e^{i (k_{\rm h}^{\rm S} L_{\rm S} + \phi)}) u^2 - (e^{i k_{\rm h}^{\rm S} L_{\rm S}} +e^{i (k_{\rm e}^{\rm S} L_{\rm S} + \phi)})v^2]}{ (e^{i k_{\rm e}^{\rm S} L_{\rm S}}-e^{i k_{\rm h}^{\rm S} L_{\rm S}})^2 u^2 v^2 - (e^{ik_{\rm h}^{\rm S} L_{\rm S}} u^2 -e^{i  k_{\rm e}^{\rm S} L_{\rm S}} v^2)^2 \,e^{i \phi}}.
\eeq
\item Transmission of an electron into NR for an incident electron from NL:
\beq\label{eq:tran-amp}
t_{\rm ee}^{\rm R}=-\frac{(u^2-v^2)^2 e^{-2i(k_{\rm e}-k_{\rm e}^{\rm S}-k_{\rm h}^{\rm S}) L_{\rm S}}e^{i\phi}} {[e^{ik_{\rm e}^{\rm S} L_{\rm S}}-e^{ik_{\rm h}^{\rm S}L_{\rm S}}]^2 u^2 v^2 - [u^2 e^{i k_{\rm h}^{\rm S}L_{\rm S}}-v^2 e^{i k_{\rm e}^{\rm S} L_{\rm S}}]^2 e^{i\phi}}.
\eeq
\end{itemize}
The amplitudes for the hole counter part can be found by taking the complex conjugate of the amplitudes for the electron part. Finally we construct the $\mathcal{S}$-matrix as
\beq
\mathcal{S} =
\begin{pmatrix}
r^{\rm L}_{\rm ee} &r^{\rm L}_{\rm eh}& t^{\prime \rm R}_{\rm ee}& t^{\prime \rm R}_{\rm eh} \\
r^{\rm L}_{\rm he} &r^{\rm L}_{\rm hh}& t^{\prime \rm R}_{he} & t^{\prime \rm R}_{\rm hh} \\
t_{\rm ee}^{\rm R} & t_{\rm eh}^{\rm R}& r_{\rm ee}^{\prime \rm L}& r_{\rm hh}^{\prime \rm L} \\
t_{\rm he}^{\rm R}  & t_{\rm hh}^{\rm R} & r_{\rm he}^{\prime \rm L}& r_{\rm hh}^{\prime \rm L}
\end{pmatrix} .
\label{smatwhole}
\eeq 
Now, due to the helicity of the edge states, we have $r^{\rm L}_{\rm ee}\!=\!r^{\rm L}_{\rm hh}\!=\!r^{\prime \rm L}_{\rm ee}\!=\!r^{\prime \rm L}_{\rm hh}\!=\!t^{\rm R}_{\rm eh}\!=\!t^{\rm L}_{\rm he}\!=\!t^{\prime \rm R}_{\rm eh}\!=\!t^{\prime \rm R}_{\rm he}\!=\!0$. After rotating the basis, we arrive at the block-diagonal form of the $\mathcal{S}$-matrix which we use to evaluate WTDs. We also rename the components of the $\mathcal{S}$-matrix here, which we have used in the main text, 
\beq
\mathcal{S} =\begin{pmatrix} t^{\prime\,R}_{ee} &r^L_{eh}&0 & 0 \\
 r_{he}^{\prime\,L}  & t^R_{hh}&  0 & 0 \\
 0& 0 & t^R_{ee}& r_{eh}^{\prime\,L} \\
0& 0 & r^L_{he} & t^{\prime\,R}_{hh}
\end{pmatrix}=\begin{pmatrix} \mathcal{S}^{\prime}_{eL,eR} &\mathcal{S}_{eL,hL}&0 & 0 \\
 \mathcal{S}_{hL,eL}^{\prime}  & \mathcal{S}_{hR,hR}&  0 & 0 \\
 0& 0 & \mathcal{S}_{eL,eR}& \mathcal{S}_{eL,hL}^{\prime} \\
0& 0 & \mathcal{S}_{hL,eL} & \mathcal{S}^{\prime}_{hL,hR}
\end{pmatrix}.
\label{smat}
\eeq 

\subsection{Waiting time distributions\label{apnd2}}

The waiting time is a fluctuating quantity, which must be described by a probability distribution. The waiting time distribution (WTD) is the conditional probability density of detecting a particle of type $\beta$ at time $t_\beta^e$, given that the last detection of a particle of type $\alpha$ occurred at the earlier time $t_\alpha^s$. Here, the types $\alpha$ and $\beta$ may refer to the out-going channel, the spin of the particle, and the particle being an electron or a hole. The WTD is denoted as $\mathcal{W}_{\alpha\rightarrow\beta}(t_\alpha^s,t_\beta^e)$. For the systems considered here, with no explicit time dependence, the WTD is a function only of the time difference, such that $\mathcal{W}_{\alpha\rightarrow\beta}(t_\alpha^s,t_\beta^e)= \mathcal{W}_{\alpha\rightarrow\beta}(\tau)$ with $\tau=t_\beta^e-t_\alpha^s$. 

%\subsubsection{Idle-time probability}

To evaluate the WTD, we proceed as in Ref.~\cite{Dasenbrook2015} and express the WTD as time-derivatives of the idle time probability. The idle time probability $\Pi(t^s_\alpha,t^{e}_\alpha)$ is the probability that no particles of type $\alpha$ are detected in the time interval $[t^s_\alpha,t^{e}_\alpha]$ by a detector at position $x_\alpha$. The idle time probability can be a function of several different particle types and associated time intervals. 

The idle time probability can be evaluated using scattering theory, leading to the determinant formula~\cite{Dasenbrook2015}
\begin{equation}
	\Pi(\{t^{s}_\gamma, t^{e}_\gamma\})= \det[\mathbbm{1} - \mathcal{Q}(\{t^{s}_{\gamma}, t^{e}_{\gamma} \})] ,
	\label{eq:ITP}
\end{equation}
where the set $\{t^{s}_\gamma, t^{e}_\gamma\}$ corresponds to all relevant particles $\gamma$ and associated time intervals. The Hermitian operator $\mathcal{Q}(\{t^{s}_\gamma, t^{e}_\gamma\})$ is a matrix in a combined energy and particle type representation. It has the block form
\begin{equation}
	[\mathcal{Q}]_{EE'} = \mathcal{S}^{\dagger}(E)\, \mathcal{K}(E-E')\,\mathcal{S}(E'),
	\label{eq:qmatrix}
\end{equation}
having omitted the time arguments. The scattering matrix $\mathcal{S}(E)$, for particles with excitation energy $E$, and the kernel $\mathcal{K}(E)$ are matrices in the space of particle types. The kernel is the block diagonal matrix
\begin{equation}
	\mathcal{K}(\{t^{s}_\gamma, t^{e}_\gamma\});E) =\bigoplus_{\gamma} K( t^{s}_\gamma, t^{e}_\gamma; E)\label{eq:K-kernels} ,
\end{equation}
given by the direct sum of kernels
\begin{equation}
	K(t^{s}_\gamma, t^{e}_\gamma; E)
	=\frac{\kappa}{\pi E} e^{-i \frac{E}{2} \left(t^{e}_\gamma+t^{s}_\gamma+\frac{2x_\gamma}{\hbar v_F}\right)}
	\sin \left[ \frac{E}{2}(t^{e}_\gamma-t^s_\gamma) \right],
	\label{eq:kernel-app}
\end{equation}
corresponding to each particle of type $\gamma$ detected at position $x_\gamma$. \Cref{eq:kernel-app} corresponds to \cref{eq:kernel} in the main text. 
We work close to the Fermi level, where the dispersion relation $E\!=\!\hbar v_Fk$ is linear and all quasi-particles propagate with the Fermi velocity~$v_F$. To implement the matrix in \cref{eq:qmatrix}, we discretize the transport window $[E_F,E_F+eV]$ in $N$ intervals of width $\kappa\!=\!eV/N$. The width $\kappa$ explicitly enters in \cref{eq:kernel-app}, and we always consider the limit $N\!\rightarrow\!\infty$, for which the transport is stationary. 

The WTD can be related to the idle time probability by realizing that time derivatives correspond to detection events~\cite{Haack2014}. When taking derivatives of operators and determinants, we use Jacobi's formula
\begin{equation}
	\frac{\mathrm{d}}{\mathrm{d}\tau} \det\{\check{A}(\tau)\} = \det\{\check{A}(\tau)\} \mathrm{Tr} \{\check{A}^{-1}(\tau) \frac{\mathrm{d}}{\mathrm{d}\tau} \check{A}(\tau)\} .
	\label{eq:Jacobi}
\end{equation}
Applied to \cref{eq:ITP}, the derivative takes the form
\begin{equation}
	\partial_{t_\alpha} \Pi(\{t^{s}_\gamma, t^{e}_\gamma\}) = - \Pi(\{t^{s}_\gamma, t^{e}_\gamma\}) \mathrm{Tr} \{ \mathcal{G}(\{t^{s}_\gamma, t^{e}_\gamma\})  \dot{\mathcal{Q}}_\alpha (t_\alpha) \} ,
	\label{eq:deriv}
\end{equation}
where we have defined 
\begin{equation}
	\mathcal{G}(\{t^{s}_\gamma, t^{e}_\gamma\})= [\mathbbm{1} - \mathcal{Q}(\{t^{s}_\gamma, t^{e}_\gamma\}) ]^{-1}
\end{equation}
and $ \dot{\mathcal{Q}}_\alpha = \partial_{t_\alpha} \mathcal{Q} $. After recasting the kernel definition, \cref{eq:kernel,eq:kernel-app}, as 
\begin{equation}
	K(t^{s}_\gamma, t^{e}_\gamma; E)
	=\frac{\kappa}{2 i \pi E} e^{-i \frac{E x_\gamma}{\hbar v_F} }
	\left( e^{-i E t^{e}_\gamma/\hbar } - e^{-i E t^{s}_\gamma / \hbar } \right),
	\label{eq:kernel2}
\end{equation}
it is straightforward to see that the derivative with respect to time $t_\alpha$ of the operator $\mathcal{Q}(\{t^{s}_\gamma, t^{e}_\gamma\}) $ depends only on time $t_\alpha$, that is, $ \dot{\mathcal{Q}}_\alpha = \dot{\mathcal{Q}}_\alpha (t_\alpha) $. Consequently, we have 
\begin{equation}
	\ddot{\mathcal{Q}}_{\alpha\beta} = \frac{\partial^2}{\partial_{t_\alpha} \partial_{t_\beta}}\mathcal{Q}(\{t^{s}_\gamma, t^{e}_\gamma\}) = \delta_{\alpha\beta} \frac{\partial^2}{\partial_{t_\beta}^2}\mathcal{Q}(\{t^{s}_\gamma, t^{e}_\gamma\}) . 
	\label{eq:dob-der}
\end{equation}

The first-passage time distributions are defined as 
\begin{equation}
	\tilde{\mathcal{F}}_{\alpha}(\tau) = \partial_{t^{s}_{\alpha}} \Pi\big|_{t^{s}_\alpha,t^s_\beta,t^e_\beta,\rightarrow 0}^{t^{e}_\beta\rightarrow \tau} = -\Pi\, \left. \mathrm{Tr}\left\{ \mathcal{G}\frac{\partial \mathcal{Q}}{\partial t^{s}_{\alpha}}\right\} \right|_{t^{s}_\alpha,t^s_\beta\rightarrow 0}^{t^e_\alpha\rightarrow 0, t^{e}_\beta\rightarrow \tau} 
	\label{eq:FPT1}
\end{equation}
and
\begin{equation}
	\mathcal{F}_{\beta}(\tau) =-\partial_{t^{e}_{\beta}}\Pi\big|_{t^{s}_\alpha,t^s_\beta,t^e_\beta,\rightarrow 0}^{t^{e}_\beta\rightarrow \tau}=\Pi\,  \left. \mathrm{Tr}\left\{\mathcal{G} \frac{\partial \mathcal{Q}}{\partial t^{e}_{\beta}}\right\} \right|_{t^{s}_\alpha,t^s_\beta\rightarrow 0}^{t^e_\alpha\rightarrow 0, t^{e}_\beta\rightarrow \tau} .
	\label{eq:FPT2}
\end{equation}
\Cref{eq:FPT1} is the conditional probability density that no particles of type $\beta$ are detected in the time span $[0,\tau]$, given that a particle of type $\alpha$ was detected at the initial time $t\!=\!0$. Similarly, \cref{eq:FPT2} concerns the time $\tau$ we have to wait until a particle of type $\alpha$ is detected, given that we start the clock at time $t\!=\!0$. 
For the single-channel stationary case, where all quantities depend only on the time difference $\tau= t^e-t^s$ and we have only one particle type, we find that these two probability distributions are the same, 
\[
\tilde{\mathcal{F}}_\alpha(\tau) = \mathcal{F}_\alpha(\tau) = -\partial_\tau \Pi(\tau) .
\]

The distribution of waiting times between particles of type $\alpha$ and particles of type $\beta$ can be expressed as~\cite{Dasenbrook2015}
\begin{equation}
	I_{\alpha} \mathcal{W}_{\alpha \rightarrow \beta}(\tau)  = -\partial_{t^{e}_\beta} \partial_{t^{s}_\alpha}\Pi(t_{\alpha}^s,t_{\alpha}^e;t_{\beta}^s,t_\beta^e\})\big|_{t^{s}_\alpha,t^s_\beta \rightarrow 0}^{t^e_\alpha \rightarrow 0, t^{e}_\beta\rightarrow \tau} \,,
	\label{eq:WTD}
\end{equation}
where $I_\alpha$ is the average particle current of type $\alpha$ particles, and the minus sign comes together with the derivative with respect to the starting time $t_\alpha^s$. In addition, after having performed the derivatives, we set the starting times to zero, i.e., $t_\alpha^s\!=\!t_\beta^s=0$, while for the end times we set $t^{e}_\alpha\!=\! 0$ and $t^{e}_\beta\!=\!\tau$.  The waiting time is then measured from the time when a particle of type $\alpha$ is detected until the later time when a particle of type $\beta$ is detected. During this waiting time, additional particles of type $\alpha$ may be detected, but not of type $\beta$.

By combining \cref{eq:WTD,eq:ITP}, we find
\begin{equation}
	I_{\alpha} \mathcal{W}_{{\alpha} \rightarrow {\beta}}(\tau)=\frac{1}{\Pi}\tilde{\mathcal{F}}_{\alpha}\mathcal{F}_{\beta}
	+ \left. \Pi \, \mathrm{Tr}\left\{ \mathcal{G} \frac{\partial \mathcal{Q}}{\partial t^{e}_{\beta}}\mathcal{G} \frac{\partial \mathcal{Q}}{\partial t^{s}_{\alpha}} \right\} \right|_{t^{s}_\alpha,t^s_\beta\rightarrow 0}^{t^e_\alpha\rightarrow 0, t^{e}_\beta\rightarrow \tau} \,,
	\label{eq:WTDfinal}
\end{equation}
having made repeatedly use of Jacobi's formula for derivatives of determinants. Finally, for evaluating \cref{eq:WTDfinal} we note that the average particle current of type $\alpha$ particles can be expressed as $I_\alpha\!=\!\mathcal{F}_{\alpha}(0)$. In combination, Eqs.~(\ref{eq:ITP}-\ref{eq:WTDfinal}) allow us to evaluate the distributions of waiting times for the superconducting systems that we consider in the main text. 
We have plotted them as a function of $\tau_i$ normalized by the mean time $\langle \tau_i \rangle$. 

\subsection{Joint waiting time distributions}

We can now introduce the joint waiting time distribution $\mathcal{W}_{\alpha \rightarrow \gamma \rightarrow \beta}(t^s_\alpha,t^m_\gamma,t^e_\beta)$. This probability distribution generalizes the waiting time distribution between particles of type $\alpha$ and $\beta$, \cref{eq:WTD,eq:WTDfinal}, to include the extra detection of a particle of type $\gamma$ at an intermediate time $t^m_\gamma$, such that $t^s_\alpha<t^m_\gamma<t^e_\beta$. 
To find the joint waiting time distribution, we introduce an auxiliary, \textit{virtual} particle channel of type $\gamma$. This new channel allows us to perform a third derivative, representing the intermediate detection event, by adding an extra kernel $K(t^s_\gamma,t^e_\gamma)$ to \cref{eq:ITP,eq:K-kernels}. Eventually, we set $t^s_\gamma=t^e_\gamma=t^m$, effectively closing this channel's contribution to the idle time probability. 

For a system with two particle types, with idle time probability $\Pi(t^s_1,t^e_1;t^s_2,t^e_2)$, the joint waiting time distribution can be defined as~\cite{Dasenbrook2015}
\begin{equation}
	I_{\alpha} \mathcal{W}_{\alpha \rightarrow \gamma \rightarrow \beta}(t^s,t^m,t^e)  = \partial_{t^{e}_\beta} \partial_{t^{e}_\gamma} \partial_{t^{s}_\alpha} \Pi( t_{\alpha}^s,t_{\alpha}^e;t_{\beta}^s,t_\beta^e;t_{\gamma}^s,t_{\gamma}^e ) 
	\big|_{t^{s}_\alpha \rightarrow t^s,t^{s}_\gamma \rightarrow t^m,t^s_\beta \rightarrow t^s}^{t^e_\alpha \rightarrow t^s, t^e_\gamma \rightarrow t^m, t^{e}_\beta\rightarrow t^e}  ,
	\label{eq:JWTD}
\end{equation}
where we have extended the idle time probability to include the auxiliary channel. From \cref{eq:dob-der}, second partial derivatives over the auxiliary channel are zero. In fact, all second derivatives are zero, since $t^s_\alpha$ and $t^e_\beta$ are independent variables even when $\alpha=\beta$. Consequently, the joint waiting time distribution has the following structure:
\begin{align}
	\partial_3 \partial_2 \partial_1 \Pi = {}& 
	-\Pi \mathrm{Tr}\{ \mathcal{G} \dot{Q}_3 \} \mathrm{Tr}\{ \mathcal{G} \dot{Q}_2 \} \mathrm{Tr}\{ \mathcal{G} \dot{Q}_1 \} 
	- \Pi \mathrm{Tr}\{ \mathcal{G} \dot{Q}_3 \mathcal{G} \dot{Q}_2 \mathcal{G} \dot{Q}_1 + \mathcal{G} \dot{Q}_3 \mathcal{G} \dot{Q}_1 \mathcal{G} \dot{Q}_2  \} \notag \\
	& + \Pi \mathrm{Tr}\{ \mathcal{G} \dot{Q}_3 \} \mathrm{Tr}\{ \mathcal{G} \dot{Q}_2 \mathcal{G} \dot{Q}_1 \} 
	+ \Pi \mathrm{Tr}\{ \mathcal{G} \dot{Q}_2 \} \mathrm{Tr}\{ \mathcal{G} \dot{Q}_3 \mathcal{G} \dot{Q}_1 \}
	+ \Pi \mathrm{Tr}\{ \mathcal{G} \dot{Q}_1 \} \mathrm{Tr}\{ \mathcal{G} \dot{Q}_3 \mathcal{G} \dot{Q}_2 \} ,
	\label{eq:jointWTD-gen}
\end{align}
where $\partial_{j}$, with $j=1,2,3$, are time derivatives over some generic and independent variables. It is important to stress that both $\Pi$ and $\mathcal{G}$ in \cref{eq:jointWTD-gen} depend only on the time coordinates of the original channels, since both $t^{s,e}_\gamma\rightarrow t^m$. That is, $\Pi=\Pi(t_{\alpha}^s,t_{\alpha}^e;t_{\beta}^s,t_\beta^e)$ and $\mathcal{G}=\mathcal{G}(t_{\alpha}^s,t_{\alpha}^e;t_{\beta}^s,t_\beta^e)$. The only dependence on $t^m$ is thus on the intermediate derivative, $\dot{Q}_2$. 

Based on the general structure of \cref{eq:jointWTD-gen}, we denote as $\gamma$ the auxiliary channel where the intermediate detection event takes place, and define the same channel joint waiting time distribution
\begin{equation}
	I_{\alpha} \mathcal{W}_{\alpha \rightarrow \gamma \rightarrow \alpha}(\tau_1,\tau_2)  = \partial_{t^{e}_\alpha} \partial_{t^{e}_\gamma} \partial_{t^{s}_\alpha} \Pi(t_{\alpha}^s,t_{\alpha}^e;t_{\beta}^s,t_\beta^e;t_{\gamma}^s,t_{\gamma}^e) 
	\big|_{t^{s}_\alpha \rightarrow 0, t^s_\gamma \rightarrow \tau_1, t^{s}_\beta \rightarrow 0}^{t^e_\alpha \rightarrow \tau_1+\tau_2, t^{e}_\gamma \rightarrow \tau_1, t^{e}_\beta \rightarrow 0}  ,
	\label{eq:JWTD11}
\end{equation}
and the different channel joint waiting time distribution
\begin{equation}
	I_{\alpha} \mathcal{W}_{\alpha \rightarrow \gamma \rightarrow \beta}(\tau_1,\tau_2)  = \partial_{t^{e}_{\beta}} \partial_{t^{e}_\gamma} \partial_{t^{s}_\alpha} \Pi(t_{\alpha}^s,t_{\alpha}^e;t_{\beta}^s,t_\beta^e;t_{\gamma}^s,t_{\gamma}^e ) 
	\big|_{t^{s}_\alpha \rightarrow 0, t^s_\gamma \rightarrow \tau_1, t^{s}_\beta \rightarrow 0}^{t^e_\alpha \rightarrow 0, t^{e}_\gamma \rightarrow \tau_1, t^{e}_\beta \rightarrow \tau_1+\tau_2}  ,
	\label{eq:JWTD12}
\end{equation}
with $\beta\neq\alpha$. 
Here, we have assumed a stationary case and set the initial time $t^s_\alpha=0$. The intermediate detection time interval is then $t^m-t^s_\alpha=\tau_1$, and the total time interval is $t^e_{\alpha,\beta}-t^s_\alpha=\tau_1+\tau_2$. 

%%%%%%%%%%%%%%%%%%%%%%%

%------------------------------------------------------
\section{\label{sec:apnd3}Scattering probabilities} 
%------------------------------------------------------

In this section, we show some results of the scattering probabilities in our interferometer to support our discussions on WTDs in the main text.
 %%%%%%%%%%%%%%%%%%%% FIGURE OF BAND STRUCTURE %%%%%%%%%%%
\begin{figure}[!thpb]
\centering
\includegraphics[scale=0.6]{figS1a.pdf}~~~
\includegraphics[scale=0.6]{figS1b.pdf}
\caption{(a) Reflection ($|r_{\rm eh}^{\rm L}|^2$) and (b) transmission ($|t_{\rm ee}^{\rm R}|^2$) probabilities as a function of the energy for $L_{\rm S}=0.2$ (green) and $1.5$ (magenta) at $\phi\!=\!0$ (dashed lines) and $\phi\!=\!\pi$ (solid lines).}
\label{fig:amplitude}
\end{figure}
%%%%%%%%%%%%%%%%%%%%%%%%%%%%%%%%%%%%%%%%%%%%%%%%%

In \cref{fig:amplitude}, we plot the energy dependence of the Andreev reflection probability $|r^{\rm L}_{\rm eh}|^2$, see \cref{eq:ref-amp}, at the left lead NL and the transmission $|t^{\rm R}_{\rm ee}|^2$, see \cref{eq:tran-amp}, at the right lead NR for $\phi\!=\!0$ and $\phi\!=\!\pi$. 
At $\phi\!=\!0$, we see that the Andreev reflection amplitude is constant (close to unity) within the superconducting gap $\Delta$ and decreases outside of the gap, followed by oscillations. The Andreev reflection for a long junction $L_{\rm s}\!\ge\!\xi$, \cref{fig:amplitude}(a), is similar to that of an ordinary normal metal-superconductor junction~\cite{Blonder1982}. Within the superconducting gap, Andreev reflection is the most likely process and thus its probability is close to unity. However, the Andreev reflection amplitude in our QSHI-based Andreev interferometer is very sensitive to the lengths of the two superconductors, $L_\text{S}$. As soon as we set $L_{\rm S}<\xi$, the probability of Andreev reflection decreases. The more we shorten the superconductor, the more the probability deviates from unity. For $L_{\rm S}\!\ll\!\xi$, it is highly flattened even beyond the gap. The mostly constant energy dependence of the probabilities for $\phi=0$ explains why their related WTDs, shown with dashed lines in \cref{fig:localWTD} of the main text, follow with little deviation the characteristic WTDs of energy-independent Poisson or Wigner-Dyson distributions. 

Interestingly, in addition to its dependence with $L_{\rm S}$, the Andreev reflection probability is also sensitive to the phase difference between the two superconductors. In the presence of a finite phase difference, the probability deviates strongly from unity around $E\!=\!0$, dropping to zero at $\phi\!=\!\pi$. This phenomenon is true for all values of $L_{\rm S}$, but the dip is particularly strong for long junctions with $L_\text{S}>\xi$. This is a result of the formation of MBS at the junction, protected by the topology of the system. The phase tunable property of the MBS at the QSHI-based normal-superconductor junction is reflected in these profiles as presented in \cref{fig:amplitude}(a). 

Due to the electron/hole separation at the edges of QSHIs, the transmission spectra provide a picture complementary to the Andreev reflection. The transmission amplitude of an electron to the right lead NR, \cref{fig:amplitude}(b), is very low throughout the entire superconductor gap region, then increases for $E$ outside the gap at $\phi\!=\!0$. Again, this is similar to an ordinary normal metal-superconductor junction~\cite{Blonder1982}. The transmission probability increases when with the lengths of the two superconductors. It becomes more dramatic when we tune the phase to $\phi=\pi$. Around $E\!=\!0$, the transmission curves rise to unity with the sharpness of the peaks controlled by $L_{\rm S}$. 

The length and phase sensitivity of Andreev reflection and electron transmission amplitudes are nicely captured by the WTDs, as discussed in the main text. Even though both normal $|t^{\rm R}_{\rm ee}|^2$ and Andreev $|r^{\rm L}_{\rm eh}|^2$ probabilities feature a strong energy dependence for $\phi=\pi$, only the distribution between electron detections, $\mathcal{W}_{ee}$, shows a behavior different than the energy-independent WTDs, cf. \cref{fig:localWTD} of the main text. In particular, $\mathcal{W}_{ee}$ approximates the distribution for a single resonance~\cite{Mi2018}. We explain this difference with $\mathcal{W}_{hh}$ as follows: for sufficiently long junctions, $L_\text{S}>\xi$, $|t^{\rm R}_{\rm ee}|^2$ presents a narrow resonance at $E\sim0$ and $|r^{\rm L}_{\rm eh}|^2$ an equally narrow dip. However, any voltage bias bigger than the resonance width ($eV=\Delta/2$ in the main text) always features almost unitary transmission for some energies close to zero and almost zero reflection probability for the same energies, and, for energies away from zero but still within the gap, $|t^{\rm R}_{\rm ee}|^2$ and $|r^{\rm L}_{\rm eh}|^2$ quickly evolve to zero and one, respectively. Consequently, as long as the voltage bias is not fine-tuned to the edge of the resonance, the WTD for electrons is more sensitive to the energy dependence since it features only a few energy intervals with high probability for $E\sim0$ and the rest are almost zero. By contrast, $\mathcal{W}_{hh}$ follows the opposite behavior with relatively more energy intervals with high probability than those with zero.

 %%%%%%%%%%%%%%%%%%%%%%%%%%%%%%%%%%%%%%%%%
\begin{figure}[t!]
\centering
\includegraphics[scale=0.6]{figS2a.pdf}~~
\includegraphics[scale=0.6]{figS2b.pdf}
	\caption{(a) $\langle \tau_{e}\rangle \mathcal{W}_{ee}$ and (b) $\langle \tau_{h}\rangle \mathcal{W}_{he}$ as functions of $\tau/\langle\tau_{e}\rangle $ at $\phi\!=\!\pi$ and $L_{\rm S}\!=\!1.5\xi$ for various energy windows $eV$. 
	In each figure, five curves from bottom to top are scaled by an integer $n \in\{1,2,3,4,5\}$ consecutively for the sake of clarity.}
\label{fig:Wee_Ls}
\end{figure}
%%%%%%%%%%%%%%%%%%%%%%%%%%%%%%%%%%%%%%%%%%


%-------------------------------------------
\section{\label{sec:bias}Bias dependence of the waiting time distributions}
%-------------------------------------------
In the main text, we have shown results for $eV=\Delta/2$ only. To better understand the effect of the voltage bias, particularly for the parameter regime where the WTDs show anomalous behaviors, we now plot local $\mathcal{W}_{ee}$ and nonlocal $\mathcal{W}_{he}$ at $\phi\!=\!\pi$ for several transport windows. 
In \cref{fig:Wee_Ls}(a), $\mathcal{W}_{ee}$ shows a plateau region for $eV=\Delta/2$, which is exactly the same result shown in the main text for the same parameter values (just scaled with respect to the axis labels). 
Increasing the bias, small oscillations appear in the WTD which give rise to the plateau feature. The frequency and amplitude of the oscillations increase with $eV$, particularly for $eV\!\ge\!\Delta$, since, in addition to the zero energy anomaly from the MBSs, the transmission also features small peaks around $E\sim\Delta$, cf. \cref{fig:amplitude}(b). The period of these oscillations is related to the energy difference between peaks, that is, it is inversely proportional to $\Delta$~\cite{Haack2014,Mi2018,Schulz2022}. 

We have already described the anomalous behavior of $\mathcal{W}_{he}$ at $\phi=\pi$ in the main text. We now plot the nonlocal WTD with the same parameters for other bias values and concentrate on short waiting times $\tau$; longer waiting times are mostly unaffected by the bias voltage~\cite{Mi2018}. 
With increasing the bias, $\mathcal{W}_{he}(0)$ increases, see Ref.~[\onlinecite{Mi2018}], and so does the value of the dip or minimum in the nonlocal WTD. We also note an oscillatory behavior for larger biases. Therefore, the behavior of the WTDs for voltages larger than the one used in the main text, but that are still comparable to the superconducting gap, does not affect our main conclusions. 

%--------------------------------------------------
\section{Nonlocal joint waiting time correlations\label{sec:joint-wtd}}
%--------------------------------------------------
Finally, we show some results for the correlations between transmitted electrons and/or reflected holes to complete \cref{WTDs_All}. 

First, the sequential tunneling of three electrons into the right lead is almost uncorrelated, i.e., $\delta\mathcal{W}_{eee}(\tau_1,\tau_2)\simeq \mathcal{W}_{ee}(\tau_1)\mathcal{W}_{ee}(\tau_2)$, as seen in \cref{fig:Weee_eV}(a) and \cref{fig:Weee_eV}(b) for $\phi=0$ and $\phi=\pi$, respectively. 

Next, we show that the correlation between waiting times for the hole-electron-hole sequence of detections, i.e., $\delta\mathcal{W}_{heh}$, is not informative for the Majorana states. 
In the absence of any phase, there is a very small correlation between two nonlocal waiting times when both of them are at very short waiting times, i.e., $\tau_{1(2)} < \langle \tau_{e(h)} \rangle$, see \cref{fig:Weee_eV}(c). 
However, when $\phi=\pi$, the correlation dies out. Mostly negative values of the correlation exist indicating dominance of the uncorrelated part $\mathcal{W}_{he}(\tau_1)\mathcal{W}_{eh}(\tau_2)$, as observed in \cref{fig:Weee_eV}(d). 
%%%%%%%%%%%%%%%%%%%%%%%%%%%%%%%%%%%%%%%%%%%
\begin{figure}[!thpb]
	\centering
	\includegraphics[scale=0.38]{figS3a.png}
	\includegraphics[scale=0.38]{figS3b.png}\\
	\includegraphics[scale=0.38]{figS3c.png}
	\includegraphics[scale=0.38]{figS3d.png}
	\caption{Correlations between waiting times for the detection of three consecutive electrons $\mathcal{W}_{\rm eee}$ and hole-electron-hole $\mathcal{W}_{\rm heh}$ for $L_{\rm S}\!=\!1.5$ as functions of $\tau_1$ and $\tau_2$ at $\phi=0$ (a,c) and $\phi=\pi$ (b,d).}
	\label{fig:Weee_eV}
\end{figure}
%%%%%%%%%%%%%%%%%%%%%%%%%%%%%%%%%%%%%%%%%%%%

%\end{appendix}

\end{document}
